\begin{proof}[Proof of \cref{obj:lem:dictionary-score-rate}]
Throughout, zeroth powers of nonnegative numbers are interpreted as \(1\). For the moment calculations, write
\[
I_s(q):=\int_{-1/2}^{1/2}|t|^s q(t)\,dt,
\qquad 0\le s\le3.
\]
For a public inverse pair, use the factor-\(16\) second-moment quantity of \cref{obj:synth_4}:
\[
V_q^{(16)}:=16\int_{-1/2}^{1/2}|t|^\kappa
E_Z[\ell_q(t+\sigma Z)^2]\,dt,
\qquad Z\sim N(0,1).
\]
The integral is understood as an extended nonnegative integral. The reference-score constants below are public and may depend on the fixed exponents \(\beta,\kappa\); the final public-score constant also depends on \(K\).

\begin{enumerate}
\item \emph{Admissibility of the dictionary pairs.}
Fix \(n\ge1\) and \(\sigma\in[0,1/4]\). We consider the three families in \cref{obj:def:weight-dictionary}.

For the constant pair, the polynomial construction in \cref{obj:synth_2} at index \(m=1\) gives
\[
q_1^{\mathrm M}(t)=\ell_1^{\mathrm M}(t)=1.
\]
The integrand defining \(I_\kappa(1)\) is continuous, nonnegative, and positive away from zero. Since the latent interval has length one and \(|t|^\kappa\le1\) there,
\[
0<I_\kappa(1)\le1,
\qquad
E_Z[|1|]=E_Z[1]=1,
\qquad
V_1^{(16)}=16I_\kappa(1)<\infty.
\]
Thus the constant pair satisfies the required conditions.

For a Fourier entry, let \(h=2^{-k}\in[n^{-2},1/4]\). Use the function \(Q\) and the pair \(q_h^{\mathrm F},\ell_h^{\mathrm F}\) from \cref{obj:synth_1}. The elementary bounds for sine imply, for \(0\le s\le3\),
\[
0\le Q(u)\le1,
\qquad
Q(u)\le |u|^{-6}\quad(u\ne0),
\qquad
0\le |u|^sQ(u)\le\frac{8}{(1+|u|)^3}.
\]
For the last bound, split at \(|u|=1\): on the central interval the product is at most one, and outside it the product is at most \(|u|^{-3}\). In particular, define
\[
c_{\mathrm F,s}:=\int_{-1/2}^{1/2}|u|^sQ(u)\,du,
\qquad
C_{\mathrm F,s}:=\int_{\mathbb R}|u|^sQ(u)\,du.
\]
Continuity and \(Q(0)=1\) give \(c_{\mathrm F,s}>0\), while the cubic envelope gives \(C_{\mathrm F,s}<\infty\). Changing variables \(t=hu\), and using \(h\le1\), yields
\[
\begin{aligned}
I_s(q_h^{\mathrm F})
&=h^{s+1}
  \int_{-1/(2h)}^{1/(2h)}|u|^sQ(u)\,du,\\
c_{\mathrm F,s}h^{s+1}
&\le I_s(q_h^{\mathrm F})
 \le C_{\mathrm F,s}h^{s+1}.
\end{aligned}
\]
These inequalities give the positive finite weighted moment.

We also record the spectral support needed below. With the transform convention of \cref{obj:synth_1}, define
\[
\varpi_h(\omega):=\frac h2\mathbf1_{[-1/h,\,1/h]}(\omega),
\qquad \omega\in\mathbb R.
\]
Direct integration gives
\[
\int_{\mathbb R}\varpi_h(\omega)e^{it\omega}\,d\omega
=\frac{\sin(t/h)}{t/h},
\]
with value one at \(t=0\). Taking six products and applying Fubini therefore gives
\[
q_h^{\mathrm F}(t)
=\int_{\mathbb R}\varpi_h^{*6}(\omega)e^{it\omega}\,d\omega,
\qquad
\operatorname{supp}\widehat q_h^{\mathrm F}
\subseteq[-6/h,6/h].
\]
Here convolution adds the support intervals; Fourier inversion gives the asserted support of the transform. Define
\[
G_{h,\sigma}(\omega)
:=\widehat q_h^{\mathrm F}(\omega)e^{\sigma^2\omega^2/2},
\qquad
\mathfrak l_{h,\sigma}(v)
:=\frac1{2\pi}\int_{\mathbb R}G_{h,\sigma}(\omega)e^{iv\omega}\,d\omega.
\]
The multiplier is continuous and compactly supported, so
\[
\ell_h^{\mathrm F}(v)=\operatorname{Re}\mathfrak l_{h,\sigma}(v),
\qquad
|\ell_h^{\mathrm F}(v)|
\le\frac1{2\pi}\int_{\mathbb R}|G_{h,\sigma}(\omega)|\,d\omega<\infty.
\]
This uniform bound proves Gaussian absolute integrability at every displacement and gives
\[
V_{q_h^{\mathrm F}}^{(16)}
\le
16I_\kappa(1)
\left(\frac1{2\pi}\int_{\mathbb R}|G_{h,\sigma}(\omega)|\,d\omega\right)^2
<\infty.
\]
The exact expectation identity is supplied by \cref{obj:synth_1}. Continuity supplies measurability of both functions.

For a polynomial entry, dictionary membership supplies a positive odd index \(m=2^j+1\). By \cref{obj:synth_2}, \(q_m^{\mathrm M}\) and \(\ell_m^{\mathrm M}\) are polynomials, and
\[
q_m^{\mathrm M}(t)\ge0,
\qquad
q_m^{\mathrm M}(0)=1.
\]
Consequently the weighted integrand is continuous on the latent interval, nonnegative, and positive on an interval away from zero. Hence
\[
0<I_\kappa(q_m^{\mathrm M})<\infty.
\]
At \(\sigma=0\), the inverse-heat sum reduces to
\[
\ell_m^{\mathrm M}(t)=q_m^{\mathrm M}(t),
\qquad
E_Z[\ell_m^{\mathrm M}(t+0Z)]=q_m^{\mathrm M}(t).
\]
For \(\sigma>0\), \cref{obj:synth_2} supplies the exact Gaussian expectation identity. In both cases, finite polynomial expansion and the Gaussian absolute moments give
\[
E_Z\!\left[|\ell_m^{\mathrm M}(t+\sigma Z)|\right]<\infty,
\qquad
16\int_{-1/2}^{1/2}|t|^\kappa
E_Z\!\left[\ell_m^{\mathrm M}(t+\sigma Z)^2\right]dt<\infty.
\]
Polynomial continuity supplies measurability. This proves pair admissibility and second-moment finiteness for every dictionary member, including when either indexed family is empty.
% lean: CausalSmith.Stat.NoisydoseWeakdesignTransition.dictionary_pair_admissible

\item \emph{Integrability under structural laws.}
Fix a dictionary pair \((q,\ell_q)\) and a law \(P\in\mathcal P_{K,\beta,\kappa,\sigma}\). Write \(X\) for the stratum coordinate, \(A\) for the latent dose, and \(Z\) for the standardized error, and define
\[
\mathcal X:=\{0,1\},
\qquad
a_0:=\frac12,
\qquad
T_{\mathrm{lat}}:=A-a_0,
\qquad
V:=A+\sigma Z-a_0.
\]
The independence and Gaussian conditions in \cref{obj:ass:error-independence,obj:ass:gaussian-channel} imply
\[
\mathcal L_P(T_{\mathrm{lat}},Z)
=\mathcal L_P(T_{\mathrm{lat}})\otimes N(0,1).
\]
For each \(x\in\mathcal X\), write \(g_x\) for the conditional Lebesgue density of \(T_{\mathrm{lat}}\) given \(X=x\), so that
\[
\mathcal L_P(T_{\mathrm{lat}}\mid X=x)(dt)=g_x(t)\,dt.
\]
The positive stratum probabilities supplied by \cref{obj:ass:stratum-overlap} make these conditional laws well defined. Since \(a_0=1/2\), \cref{obj:ass:latent-support} gives
\[
P\{T_{\mathrm{lat}}\in[-1/2,1/2]\}=1.
\]
Thus, for any nonnegative measurable function \(f\), conditioning on the stratum gives the following equality with the integral restricted to the latent interval; the upper density envelope in \cref{obj:ass:weak-design}, with \(p_x:=P(X=x)\le1\), gives the inequality:
\[
\begin{aligned}
E_P[\mathbf1\{X=x\}f(T_{\mathrm{lat}})]
&=p_x\int_{-1/2}^{1/2}g_x(t)f(t)\,dt\\
&\le c_{\mathrm{hi}}\int_{-1/2}^{1/2}|t|^\kappa f(t)\,dt.
\end{aligned}
\]
Summing the two strata and applying this inequality to
\[
f_q(t):=E_Z[\ell_q(t+\sigma Z)^2],
\qquad t\in\mathbb R,
\]
proves, by the product-law identity and Tonelli,
\[
E_P[\ell_q(V)^2]
=E_P[f_q(T_{\mathrm{lat}})]
\le\frac{c_{\mathrm{hi}}}{16}
\bigl(V_q^{(16)}+V_q^{(16)}\bigr)<\infty.
\]
The compositions involved are measurable. The inclusion \(L^2(P)\subseteq L^1(P)\) on a probability space now gives
\[
E_P[|\ell_q(V)|]
\le\sqrt{E_P[\ell_q(V)^2]}
\le\sqrt{\frac{c_{\mathrm{hi}}}{8}V_q^{(16)}}<\infty.
\]
This establishes the law-integrability assertion.
% lean: CausalSmith.Stat.NoisydoseWeakdesignTransition.dictionary_pair_law_admissible

\item \emph{The Fourier score bound.}
For finite-moment pairs and \(n\ge1\), define the reference quantities used to tune the dictionary:
\[
b_q^{\mathrm{ref}}:=\frac1{16}I_\kappa(q),
\qquad
B_q^{\mathrm{ref}}:=64\frac{I_{\kappa+\beta}(q)}{I_\kappa(q)},
\qquad
A_q^{\mathrm{ref}}:=B_q^{\mathrm{ref}}
+\frac{12}{b_q^{\mathrm{ref}}}\sqrt{\frac{V_q^{(16)}}{n}}
+2\sqrt{\frac3n}.
\]
Since \(0\le\kappa+\beta\le3\), the Fourier moment bounds already proved give
\[
B_{q_h^{\mathrm F}}^{\mathrm{ref}}
\le C_{\mathrm F,bias}h^\beta,
\qquad
C_{\mathrm F,bias}:=
\frac{64C_{\mathrm F,\kappa+\beta}}{c_{\mathrm F,\kappa}}>0.
\]

To control the variance, define the spectral profile constants
\[
a_{\mathrm F}:=\int_{\mathbb R}Q(u)\,du,
\qquad
b_{\mathrm F}:=\int_{\mathbb R}|u|Q(u)\,du.
\]
Differentiation under the Fourier integral, justified by the absolute moments through order three, gives
\[
\widehat q_h^{\mathrm F}\in C^3(\mathbb R),
\qquad
|\widehat q_h^{\mathrm F}(\omega)|\le a_{\mathrm F}h,
\qquad
|(\widehat q_h^{\mathrm F})'(\omega)|\le b_{\mathrm F}h^2.
\]
Both \(G_{h,\sigma}\) and its derivatives have compact support. Twice integrating by parts gives, for \(v\ne0\),
\[
|\mathfrak l_{h,\sigma}(v)|
\le\frac1{2\pi}
\min\left\{
\int_{\mathbb R}|G_{h,\sigma}(\omega)|\,d\omega,\,
\frac1{v^2}\int_{\mathbb R}|G_{h,\sigma}''(\omega)|\,d\omega
\right\}.
\]
Together with the direct bound at \(v=0\), this proves absolute integrability of the inverse. The same argument applied to \(G_{h,\sigma}'\) proves absolute integrability of its inverse.

For \(0\le s\le2\), define the physical-space energies
\[
\mathfrak e_{s,h,\sigma}
:=\int_{\mathbb R}|v|^s|\mathfrak l_{h,\sigma}(v)|^2\,dv.
\]
Parseval's identity and the inverse-transform derivative identity give
\[
\mathfrak e_{0,h,\sigma}
=\frac1{2\pi}\int_{\mathbb R}|G_{h,\sigma}(\omega)|^2\,d\omega,
\qquad
\mathfrak e_{2,h,\sigma}
=\frac1{2\pi}\int_{\mathbb R}|G_{h,\sigma}'(\omega)|^2\,d\omega.
\]
On the spectral support, \(|\omega|\le6/h\). Thus
\[
|G_{h,\sigma}(\omega)|
\le a_{\mathrm F}h\,e^{18(\sigma/h)^2}.
\]
Also,
\[
G_{h,\sigma}'
=e^{\sigma^2\omega^2/2}
\left((\widehat q_h^{\mathrm F})'
      +\sigma^2\omega\widehat q_h^{\mathrm F}\right),
\]
so
\[
\begin{aligned}
|G_{h,\sigma}'(\omega)|
&\le h^2\left(b_{\mathrm F}
             +6a_{\mathrm F}(\sigma/h)^2\right)e^{18(\sigma/h)^2}\\
&\le h^2(b_{\mathrm F}+6a_{\mathrm F})
       (1+\sigma/h)^2e^{18(\sigma/h)^2}.
\end{aligned}
\]
Integrating these squared bounds over an interval of length \(12/h\) yields
\[
\begin{aligned}
\mathfrak e_{0,h,\sigma}
&\le\frac6\pi a_{\mathrm F}^2h\,e^{36(\sigma/h)^2},\\
\mathfrak e_{2,h,\sigma}
&\le\frac6\pi(b_{\mathrm F}+6a_{\mathrm F})^2
h^3(1+\sigma/h)^4e^{36(\sigma/h)^2}.
\end{aligned}
\]
For \(0\le\kappa\le2\), splitting at \(|v|=h\) gives
\[
|v|^\kappa\le h^\kappa+h^{\kappa-2}v^2.
\]
Consequently,
\[
\begin{aligned}
\mathfrak e_{\kappa,h,\sigma}
&\le h^\kappa\mathfrak e_{0,h,\sigma}
    +h^{\kappa-2}\mathfrak e_{2,h,\sigma}\\
&\le\frac6\pi
\left(a_{\mathrm F}^2+(b_{\mathrm F}+6a_{\mathrm F})^2\right)
h^{\kappa+1}(1+\sigma/h)^4e^{36(\sigma/h)^2}.
\end{aligned}
\]

The power inequality
\[
|v-\sigma z|^\kappa
\le2\left(|v|^\kappa+\sigma^\kappa|z|^\kappa\right)
\]
holds throughout the stated range, including \(\kappa=0\). Moreover,
\[
E_Z[|Z|^\kappa]\le E_Z[1+Z^2]=2.
\]
These inequalities also supply joint integrability after translation: the dominating terms are integrable products of a Gaussian moment and one of the finite energies above. Extending the latent integral to the real line, translating \(v=t+\sigma z\), and using
\(|\operatorname{Re}\mathfrak l|\le|\mathfrak l|\), we obtain
\[
\begin{aligned}
V_{q_h^{\mathrm F}}^{(16)}
&\le16E_Z\int_{\mathbb R}|v-\sigma Z|^\kappa
                      |\mathfrak l_{h,\sigma}(v)|^2\,dv\\
&\le32\left(\mathfrak e_{\kappa,h,\sigma}
                 +2\sigma^\kappa\mathfrak e_{0,h,\sigma}\right).
\end{aligned}
\]
Since
\[
\sigma^\kappa h
=h^{\kappa+1}(\sigma/h)^\kappa,
\qquad
(\sigma/h)^\kappa\le(1+\sigma/h)^2,
\]
enlarging the powers \(2\) and \(4\) to \(10\) gives
\[
V_{q_h^{\mathrm F}}^{(16)}
\le C_{\mathrm F,var}\,
h^{\kappa+1}(1+\sigma/h)^{10}e^{36(\sigma/h)^2},
\]
where
\[
C_{\mathrm F,var}
:=\frac{192}{\pi}
\left(3a_{\mathrm F}^2+(b_{\mathrm F}+6a_{\mathrm F})^2\right)+1>0.
\]
% lean: CausalSmith.Stat.NoisydoseWeakdesignTransition.ellF_variance_bound

The denominator moment bound implies
\[
\frac{12}{b_{q_h^{\mathrm F}}^{\mathrm{ref}}}
\le\frac{192}{c_{\mathrm F,\kappa}h^{\kappa+1}}.
\]
Define, for \(n\ge1\), \(h>0\), and \(\sigma\ge0\),
\[
\mathfrak s_{n,\sigma}(h)
:=\sqrt{\frac{h^{-\kappa-1}(1+\sigma/h)^{10}
                         e^{36(\sigma/h)^2}}{n}},
\]
and set
\[
C_{\mathrm F,score}
:=C_{\mathrm F,bias}
+\frac{192\sqrt{C_{\mathrm F,var}}}{c_{\mathrm F,\kappa}}
+2\sqrt3+1.
\]
Combining the bias, denominator, and variance bounds in the reference score gives
\[
A_{q_h^{\mathrm F}}^{\mathrm{ref}}
\le C_{\mathrm F,score}
\left(h^\beta+\mathfrak s_{n,\sigma}(h)+n^{-1/2}\right),
\qquad
0<h\le\frac14,\quad 0\le\sigma\le\frac14.
\]
% lean: CausalSmith.Stat.NoisydoseWeakdesignTransition.qF_certificate_bound

\item \emph{The direct Fourier witness.}
Define
\[
d:=2\beta+\kappa+1,
\qquad
D_n:=n^{-1/d},
\qquad n\ge1.
\]
The parameter ranges imply \(1<d\le5\). The direct scale satisfies
\[
D_n^d=\frac1n,
\qquad
\frac{D_n^{-\kappa-1}}n=D_n^{2\beta},
\qquad
n^{-1/2}\le D_n^\beta.
\]
The last inequality follows from \(D_n\le1\) and \(d\ge2\beta\).

For sufficiently large \(n\),
\[
0<D_n,
\qquad
n^{-2}\le D_n\le\frac18;
\]
the lower bound follows by comparing exponents, and the upper bound follows from \(D_n\to0\). Choose a dyadic bandwidth by
\[
2^{-(k+1)}<D_n\le2^{-k},
\qquad
h:=2^{-k}.
\]
Then
\[
D_n\le h<2D_n,
\qquad
n^{-2}\le h\le\frac14,
\]
so the pair belongs to the dictionary by \cref{obj:def:weight-dictionary}.

Suppose \(\sigma\le D_n\). For fixed nonnegative \(\sigma\), every factor inside \(\mathfrak s_{n,\sigma}(h)^2\) decreases as \(h\) increases. Hence
\[
\begin{aligned}
h^\beta&\le2^\beta D_n^\beta,\\
\mathfrak s_{n,\sigma}(h)
&\le\mathfrak s_{n,\sigma}(D_n)
 \le\sqrt{2^{10}e^{36}}\,D_n^\beta,
\end{aligned}
\]
where the second bound uses \(0\le\sigma/D_n\le1\) and the variance balance above. Therefore, with
\[
C_{\mathrm{direct}}
:=C_{\mathrm F,score}
\left(2^\beta+\sqrt{2^{10}e^{36}}+1\right)>0,
\]
the rounded Fourier entry satisfies
\[
A_{q_h^{\mathrm F}}^{\mathrm{ref}}
\le C_{\mathrm{direct}}D_n^\beta
=C_{\mathrm{direct}}\rho_{n,\sigma}.
\]
The equality is the direct branch of \cref{obj:def:frontier-rate}. Thus there is an integer \(n_{\mathrm{direct}}\), independent of \(\sigma\), beyond which this witness is available whenever \(\sigma\le D_n\).
% lean: CausalSmith.Stat.NoisydoseWeakdesignTransition.direct_fourier_dictionary_witness

\item \emph{The intermediate Fourier witness.}
For \(n\ge1\) and \(0<\sigma\le1/4\), define
\[
N_{\mathrm F}:=n\sigma^d,
\qquad
\Lambda_{\mathrm F}:=\log(e+N_{\mathrm F}),
\qquad
F_n(\sigma):=\frac{\sigma}{\sqrt{\Lambda_{\mathrm F}}}.
\]
When \(D_n<\sigma\), the identity \(nD_n^d=1\) implies \(N_{\mathrm F}\ge1\), and hence \(\Lambda_{\mathrm F}\ge1\).

We first establish the uniform logarithmic bound used for tuning. Define
\[
C_{\log}:=1024\cdot8!\left(\frac43\right)^8(e+1)>0.
\]
For \(N_{\mathrm F}\ge1\), the inequalities \(d\le5\) and \(\Lambda_{\mathrm F}\ge1\) give
\[
\Lambda_{\mathrm F}^{d/2}\le\Lambda_{\mathrm F}^3,
\qquad
\left(1+\frac{\sqrt{\Lambda_{\mathrm F}}}{12}\right)^{10}
\le1024\Lambda_{\mathrm F}^5.
\]
The eighth term of the exponential series gives
\[
\Lambda_{\mathrm F}^8
\le8!\left(\frac43\right)^8e^{3\Lambda_{\mathrm F}/4}.
\]
Using \(e^{\Lambda_{\mathrm F}}=e+N_{\mathrm F}\) and
\(e+N_{\mathrm F}\le(e+1)N_{\mathrm F}\), we conclude
\[
\frac{\Lambda_{\mathrm F}^{d/2}
\left(1+\sqrt{\Lambda_{\mathrm F}}/12\right)^{10}
e^{\Lambda_{\mathrm F}/4}}{N_{\mathrm F}}
\le C_{\log}.
\]
% lean: CausalSmith.Stat.NoisydoseWeakdesignTransition.fourier_log_tuning_bound

At the target bandwidth \(12F_n(\sigma)\),
\[
\frac{\sigma}{12F_n(\sigma)}
=\frac{\sqrt{\Lambda_{\mathrm F}}}{12},
\qquad
36\left(\frac{\sigma}{12F_n(\sigma)}\right)^2
=\frac{\Lambda_{\mathrm F}}4,
\qquad
\frac{F_n(\sigma)^{-d}}n
=\frac{\Lambda_{\mathrm F}^{d/2}}{N_{\mathrm F}}.
\]
Since \(12F_n(\sigma)\ge F_n(\sigma)\) and \(d>0\), these identities yield
\[
\frac{\mathfrak s_{n,\sigma}(12F_n(\sigma))^2}
     {(12F_n(\sigma))^{2\beta}}
\le
\frac{\Lambda_{\mathrm F}^{d/2}
\left(1+\sqrt{\Lambda_{\mathrm F}}/12\right)^{10}
e^{\Lambda_{\mathrm F}/4}}{N_{\mathrm F}}
\le C_{\log}.
\]
Thus
\[
\mathfrak s_{n,\sigma}(12F_n(\sigma))
\le\sqrt{C_{\log}}\,(12F_n(\sigma))^\beta.
\]
Dropping the two factors that are at least one from the logarithmic bound also gives
\[
\frac{F_n(\sigma)^{-d}}n\le C_{\log}.
\]
Because \(0<F_n(\sigma)\le\sigma\le1/4\) and \(d\ge2\beta\), it follows that
\[
\frac1n\le C_{\log}F_n(\sigma)^d
\le C_{\log}F_n(\sigma)^{2\beta},
\qquad
n^{-1/2}\le\sqrt{C_{\log}}F_n(\sigma)^\beta.
\]

We next place the target bandwidth in the dictionary window. If
\(F_n(\sigma)\le n^{-2}\), then the inverse-power bound would imply
\[
n^{2d-1}\le C_{\log}.
\]
For \(n\ge1\), \(d>1\) gives \(n\le n^{2d-1}\). Thus, once \(n>C_{\log}\),
\[
n^{-2}\le12F_n(\sigma).
\]
If also \(\sigma\le L_n^{-1/2}\), then
\[
12F_n(\sigma)\le12\sigma\le12L_n^{-1/2}\le\frac18
\]
for sufficiently large \(n\), since \(L_n\to\infty\). These thresholds are independent of \(\sigma\).

Choose \(k\) and \(h\) by
\[
2^{-(k+1)}<12F_n(\sigma)\le2^{-k},
\qquad
h:=2^{-k}.
\]
Then
\[
12F_n(\sigma)\le h<24F_n(\sigma),
\qquad
n^{-2}\le h\le\frac14,
\]
so the pair belongs to the dictionary. Monotonicity of the stochastic factor and the Fourier reference-score bound give
\[
A_{q_h^{\mathrm F}}^{\mathrm{ref}}
\le C_{\mathrm{intermediate}}F_n(\sigma)^\beta,
\]
where
\[
C_{\mathrm{intermediate}}
:=C_{\mathrm F,score}
\left(24^\beta+\sqrt{C_{\log}}\,12^\beta+\sqrt{C_{\log}}\right)>0.
\]
Under \(D_n<\sigma\le L_n^{-1/2}\), \cref{obj:def:frontier-rate} identifies \(F_n(\sigma)^\beta=\rho_{n,\sigma}\). This proves the intermediate witness for all \(n\ge n_{\mathrm{intermediate}}\), with a threshold independent of \(\sigma\).
% lean: CausalSmith.Stat.NoisydoseWeakdesignTransition.intermediate_fourier_dictionary_witness

\item \emph{Combining the Fourier regimes.}
Set
\[
C_{\mathrm F}:=\max\{C_{\mathrm{direct}},C_{\mathrm{intermediate}}\},
\qquad
n_{\mathrm F}:=\max\{n_{\mathrm{direct}},n_{\mathrm{intermediate}}\}.
\]
For \(n\ge n_{\mathrm F}\) and \(\sigma\le L_n^{-1/2}\), split according to \(\sigma\le D_n\) or \(D_n<\sigma\). The preceding two steps supply, respectively, a dictionary Fourier pair satisfying
\[
A_{q_{2^{-k}}^{\mathrm F}}^{\mathrm{ref}}
\le C_{\mathrm F}\rho_{n,\sigma}.
\]
The enlargement of the constants is valid because \(\rho_{n,\sigma}\ge0\). The direct case includes \(\sigma=0\).
% lean: CausalSmith.Stat.NoisydoseWeakdesignTransition.fourier_dictionary_witness

\item \emph{The polynomial certificate.}
We derive the inverse-degree score bound before choosing the dictionary index. For any positive odd \(m\), the inequality
\[
|\sin(mz)|\le m|\sin z|
\]
follows by induction from the sine addition formula. Combining it with
\(|\sin(mz)|\le1\) and the representation in \cref{obj:synth_2} gives
\[
0\le q_m^{\mathrm M}(t)\le1,
\qquad
q_m^{\mathrm M}(t)\le(2m|t|)^{-6}\quad(t\ne0),
\qquad |t|\le\frac12.
\]
Jordan's inequality and \(\sin(\arcsin(2t))=2t\) imply
\[
2|t|\le|\arcsin(2t)|\le\pi|t|.
\]
Therefore, if \(|t|\le1/(2m)\), then
\[
|m\arcsin(2t)|\le\frac\pi2,
\qquad
|\sin(m\arcsin(2t))|
\ge\frac2\pi m|\arcsin(2t)|
\ge\frac4\pi m|t|,
\]
and consequently
\[
q_m^{\mathrm M}(t)\ge\left(\frac2\pi\right)^6.
\]
At \(t=0\), this follows from \(q_m^{\mathrm M}(0)=1\).

For \(0\le s\le3\), define
\[
c_{\mathrm M,s}:=\left(\frac2\pi\right)^6I_s(1)>0,
\qquad
C_{\mathrm{tail}}:=\int_{\mathbb R}\frac8{(1+|u|)^3}\,du<\infty.
\]
Integrating the central lower bound and rescaling gives
\[
I_s(q_m^{\mathrm M})\ge c_{\mathrm M,s}m^{-(s+1)}.
\]
For the upper bound, put \(u=mt\). Splitting at \(|u|=1\), the preceding upper bounds give
\[
|t|^s q_m^{\mathrm M}(t)
\le m^{-s}\frac8{(1+|mt|)^3}.
\]
Extending this nonnegative envelope to the real line and changing variables yields
\[
c_{\mathrm M,s}m^{-(s+1)}
\le I_s(q_m^{\mathrm M})
\le C_{\mathrm{tail}}m^{-(s+1)}.
\]
In particular, define
\[
C_{\mathrm M,bias}:=\frac{64C_{\mathrm{tail}}}{c_{\mathrm M,\kappa}},
\qquad
C_{\mathrm M,ratio}:=\frac{192}{c_{\mathrm M,\kappa}}.
\]
Then
\[
B_{q_m^{\mathrm M}}^{\mathrm{ref}}\le C_{\mathrm M,bias}m^{-\beta},
\qquad
\frac{12}{b_{q_m^{\mathrm M}}^{\mathrm{ref}}}
\le C_{\mathrm M,ratio}m^{\kappa+1}.
\]

For the variance estimate, define the degree bound
\[
D_{\mathrm{pol},m}:=6(m-1).
\]
We first bound the coefficients of the declared polynomial. Let the second-kind Chebyshev polynomials be defined by
\[
\mathsf U^{\mathrm{Ch}}_0(s):=1,\qquad
\mathsf U^{\mathrm{Ch}}_1(s):=2s,\qquad
\mathsf U^{\mathrm{Ch}}_{i+2}(s)
:=2s\mathsf U^{\mathrm{Ch}}_{i+1}(s)-\mathsf U^{\mathrm{Ch}}_i(s).
\]
For a polynomial \(p\), define its coefficient mass by
\[
\|p\|_{\mathrm{coef}}
:=\sum_{i\ge0}|[t^i]p(t)|;
\]
the sum is finite. The recurrence, together with the triangle inequality and submultiplicativity of coefficient mass, proves inductively that
\[
\|\mathsf U^{\mathrm{Ch}}_i\|_{\mathrm{coef}}\le3^i.
\]
Indeed, the initial masses are \(1,2\), and
\[
2\cdot3^{i+1}+3^i\le3^{i+2}.
\]
Define the even coefficients by
\[
u^{\mathrm{Ch}}_{m,i}:=[s^{2i}]\mathsf U^{\mathrm{Ch}}_{m-1}(s),
\qquad 0\le i\le m-1.
\]
The identity
\(\mathsf U^{\mathrm{Ch}}_{m-1}(\cos z)=\sin(mz)/\sin z\),
obtained from the same recurrence, gives the polynomial representation
\[
q_m^{\mathrm M}(t)
=\left[
\frac1m\sum_{i=0}^{m-1}
u^{\mathrm{Ch}}_{m,i}(1-4t^2)^i
\right]^6.
\]
Since \(\|1-4t^2\|_{\mathrm{coef}}=5\), every summand inside the brackets has coefficient mass at most
\[
\frac{3^{m-1}5^i}{m}\le\frac{15^{m-1}}m.
\]
Summing the \(m\) terms and taking the sixth power gives
\[
\|q_m^{\mathrm M}\|_{\mathrm{coef}}
\le15^{D_{\mathrm{pol},m}},
\qquad
\deg q_m^{\mathrm M}\le D_{\mathrm{pol},m}.
\]
% lean: CausalSmith.Stat.NoisydoseWeakdesignTransition.qMPoly_coeffOneNorm_le

Define the probabilists' Hermite polynomials by
\[
\mathsf{He}_i(z)
:=i!\sum_{a=0}^{\lfloor i/2\rfloor}
\frac{(-1)^a z^{i-2a}}{2^a a!(i-2a)!},
\qquad i\ge0,\ z\in\mathbb R.
\]
The finite coefficient formula gives the recurrence
\[
\mathsf{He}_0(z)=1,\qquad
\mathsf{He}_1(z)=z,\qquad
\mathsf{He}_{i+1}(z)=z\mathsf{He}_i(z)-\mathsf{He}_i'(z),
\qquad i\ge0,\ z\in\mathbb R.
\]
Differentiating this recurrence proves by induction that
\[
\mathsf{He}_0'(z)=0,\qquad
\mathsf{He}_{i+1}'(z)=(i+1)\mathsf{He}_i(z),
\qquad i\ge0.
\]
Indeed, the initial identity is \(\mathsf{He}_1'=1\); for \(i\ge1\), the preceding derivative identity gives
\[
\begin{aligned}
\mathsf{He}_{i+1}'(z)
&=\mathsf{He}_i(z)+z\mathsf{He}_i'(z)-\mathsf{He}_i''(z)\\
&=\mathsf{He}_i(z)
+i\bigl(z\mathsf{He}_{i-1}(z)-\mathsf{He}_{i-1}'(z)\bigr)\\
&=(i+1)\mathsf{He}_i(z).
\end{aligned}
\]
% lean: Causalean.Mathlib.Probability.derivative_hermite_succ

For every real polynomial \(p\), Gaussian decay gives
\[
\lim_{z\to\pm\infty}p(z)e^{-z^2/2}=0.
\]
Integration by parts, followed by Gaussian normalization, therefore yields
\[
\begin{aligned}
\int_{\mathbb R}zp(z)e^{-z^2/2}\,dz
&=-\int_{\mathbb R}p(z)\frac{d}{dz}e^{-z^2/2}\,dz\\
&=\int_{\mathbb R}p'(z)e^{-z^2/2}\,dz,\\
E_Z[Zp(Z)]&=E_Z[p'(Z)].
\end{aligned}
\]
All polynomial factors have finite Gaussian absolute moments. Applying this identity to the Hermite recurrence gives
\[
E_Z[\mathsf{He}_0(Z)]=1,\qquad
E_Z[\mathsf{He}_{i+1}(Z)]
=E_Z[Z\mathsf{He}_i(Z)]-E_Z[\mathsf{He}_i'(Z)]=0,
\qquad i\ge0.
\]
% lean: Causalean.Mathlib.Probability.integral_hermite_succ

Define the Gaussian inner products by
\[
\mathfrak h_{i,a}:=E_Z[\mathsf{He}_i(Z)\mathsf{He}_a(Z)],
\qquad i,a\in\mathbb Z_{\ge0}.
\]
Applying integration by parts to the polynomial product
\(\mathsf{He}_i\mathsf{He}_{a+1}\), the recurrence and derivative identity give
\[
\begin{aligned}
\mathfrak h_{i+1,a+1}
&=E_Z[Z\mathsf{He}_i(Z)\mathsf{He}_{a+1}(Z)]
  -E_Z[\mathsf{He}_i'(Z)\mathsf{He}_{a+1}(Z)]\\
&=E_Z[\mathsf{He}_i'(Z)\mathsf{He}_{a+1}(Z)
       +\mathsf{He}_i(Z)\mathsf{He}_{a+1}'(Z)]
  -E_Z[\mathsf{He}_i'(Z)\mathsf{He}_{a+1}(Z)]\\
&=(a+1)\mathfrak h_{i,a}.
\end{aligned}
\]
% lean: Causalean.Mathlib.Probability.integral_hermite_succ_mul_succ
The mean identities supply the boundary values
\[
\mathfrak h_{0,0}=1,\qquad
\mathfrak h_{i+1,0}=\mathfrak h_{0,i+1}=0,
\qquad i\ge0.
\]
Induct on the first index, allowing every second index. The boundary values handle either index being zero. When both indices are positive, the recurrence reduces both by one. On the diagonal it gives
\[
\mathfrak h_{i+1,i+1}=(i+1)\mathfrak h_{i,i}
=(i+1)i!=(i+1)!;
\]
for \(i\ne a\), the induction hypothesis gives
\[
\mathfrak h_{i+1,a+1}=(a+1)\mathfrak h_{i,a}=0.
\]
Thus their Gaussian orthogonality is
\[
E_Z[\mathsf{He}_i(Z)\mathsf{He}_a(Z)]
=\mathbf1\{i=a\}\,i!.
\]
% lean: Causalean.Mathlib.Probability.integral_hermite_mul
Taylor expansion gives
\[
q_m^{\mathrm M}(t+\sigma z)
=\sum_{i=0}^{D_{\mathrm{pol},m}}
\frac{\sigma^i(q_m^{\mathrm M})^{(i)}(t)}{i!}z^i.
\]
Applying the finite inverse-heat sum from \cref{obj:synth_2} to this expansion replaces \(z^i\) by \(\mathsf{He}_i(z)\): this is the displayed finite formula for \(\mathsf{He}_i\), with the chain rule supplying the powers of \(\sigma\). Thus
\[
\ell_m^{\mathrm M}(t+\sigma Z)
=\sum_{i=0}^{D_{\mathrm{pol},m}}
\frac{\sigma^i(q_m^{\mathrm M})^{(i)}(t)}{i!}\mathsf{He}_i(Z).
\]
Squaring and using orthogonality yields
\[
E_Z[\ell_m^{\mathrm M}(t+\sigma Z)^2]
=\sum_{i=0}^{D_{\mathrm{pol},m}}
\frac{\sigma^{2i}\bigl((q_m^{\mathrm M})^{(i)}(t)\bigr)^2}{i!}.
\]
These finite identities include \(\sigma=0\).
% lean: CausalSmith.Stat.NoisydoseWeakdesignTransition.ellM_gaussian_second_moment

For \(0\le i\le D_{\mathrm{pol},m}\), write
\[
(D_{\mathrm{pol},m})_i
:=\frac{D_{\mathrm{pol},m}!}{(D_{\mathrm{pol},m}-i)!}.
\]
Differentiating the monomial expansion and using \(|t|\le1\) gives
\[
|(q_m^{\mathrm M})^{(i)}(t)|
\le\|q_m^{\mathrm M}\|_{\mathrm{coef}}(D_{\mathrm{pol},m})_i
\le15^{D_{\mathrm{pol},m}}(D_{\mathrm{pol},m})_i.
\]
Moreover,
\[
\frac{(D_{\mathrm{pol},m})_i^2}{i!}
=\binom{D_{\mathrm{pol},m}}i(D_{\mathrm{pol},m})_i
\le\binom{D_{\mathrm{pol},m}}iD_{\mathrm{pol},m}^i.
\]
The binomial theorem and \(I_\kappa(1)\le1\) therefore give
\[
\begin{aligned}
V_{q_m^{\mathrm M}}^{(16)}
&\le16I_\kappa(1)\,15^{2D_{\mathrm{pol},m}}
       (1+D_{\mathrm{pol},m}\sigma^2)^{D_{\mathrm{pol},m}}\\
&\le16\,15^{2D_{\mathrm{pol},m}}
       (1+D_{\mathrm{pol},m}\sigma^2)^{D_{\mathrm{pol},m}}.
\end{aligned}
\]
The case \(m=1\) has \(D_{\mathrm{pol},m}=0\), and the same formulas reduce to the constant pair.
% lean: CausalSmith.Stat.NoisydoseWeakdesignTransition.ellM_Vq_explicit_bound

For \(n\ge1\) and \(\sigma\in[0,1/4]\), define
\[
\Lambda_{\mathrm M,n,\sigma}:=\log(e+\sigma^2L_n).
\]
Here \(L_n\ge1\) and \(\Lambda_{\mathrm M,n,\sigma}\ge1\). Suppose
\[
m\le\frac{L_n}{4096\Lambda_{\mathrm M,n,\sigma}}.
\]
Then
\[
D_{\mathrm{pol},m}
\le\frac{6L_n}{4096\Lambda_{\mathrm M,n,\sigma}},
\qquad
1+D_{\mathrm{pol},m}\sigma^2
\le6(e+\sigma^2L_n),
\]
so
\[
\log(1+D_{\mathrm{pol},m}\sigma^2)
\le\Lambda_{\mathrm M,n,\sigma}+\log6.
\]
Using \(\log15\le15\), \(\log6\le6\), and
\(\Lambda_{\mathrm M,n,\sigma}\ge1\), we obtain the explicit calibration
\[
\begin{aligned}
&D_{\mathrm{pol},m}\log15
+\frac{D_{\mathrm{pol},m}}2
 \log(1+D_{\mathrm{pol},m}\sigma^2)\\
&\quad\le
\frac{6L_n}{4096\Lambda_{\mathrm M,n,\sigma}}
\left(15+\frac{\Lambda_{\mathrm M,n,\sigma}+6}{2}\right)
\le\frac{L_n}{16}.
\end{aligned}
\]
The last numerical inequality follows from
\(6(18+\Lambda_{\mathrm M,n,\sigma}/2)
\le256\Lambda_{\mathrm M,n,\sigma}\).
Exponentiating twice the bound gives
\[
15^{2D_{\mathrm{pol},m}}
(1+D_{\mathrm{pol},m}\sigma^2)^{D_{\mathrm{pol},m}}
\le e^{L_n/8}\le n^{1/4}
\]
whenever \(L_n=1+\log n\ge2\). Consequently,
\[
V_{q_m^{\mathrm M}}^{(16)}\le16n^{1/4},
\qquad
\sqrt{\frac{V_{q_m^{\mathrm M}}^{(16)}}{n}}\le4n^{-3/8}.
\]

Finally, the standard logarithmic-power comparison gives
\[
\frac{L_n^4}{n^{3/8}}\longrightarrow0,
\qquad
L_n^4\le n^{3/8}
\quad\text{for all sufficiently large }n.
\]
The degree condition implies \(1\le m\le L_n\). Since
\(\kappa+\beta+1\le4\), it follows that
\[
m^{\kappa+\beta+1}\le m^4\le L_n^4\le n^{3/8},
\qquad
n^{-3/8}m^{\kappa+1}\le m^{-\beta}.
\]
Also,
\[
n^{-1/2}\le n^{-3/8}
\le n^{-3/8}m^{\kappa+1}\le m^{-\beta}.
\]
Substitution into the reference score proves
\[
A_{q_m^{\mathrm M}}^{\mathrm{ref}}
\le C_{\mathrm{poly}}m^{-\beta},
\qquad
C_{\mathrm{poly}}
:=C_{\mathrm M,bias}+4C_{\mathrm M,ratio}+2\sqrt3>0,
\]
uniformly for all sufficiently large \(n\), all public \(\sigma\), and all positive odd \(m\) satisfying the degree condition.
% lean: CausalSmith.Stat.NoisydoseWeakdesignTransition.qM_tuned_certificate_bound

\item \emph{The polynomial dictionary witness.}
Define the real target index
\[
m_{\mathrm{tar},n,\sigma}
:=\frac{L_n}{4096\Lambda_{\mathrm M,n,\sigma}}.
\]
Uniformly over \(\sigma\in[0,1/4]\),
\[
1\le\Lambda_{\mathrm M,n,\sigma}\le\log(e+L_n).
\]
Because \(\log(e+L_n)/L_n\to0\), eventually
\[
\log(e+L_n)\le\frac{L_n}{16384}.
\]
For sufficiently large \(n\), the inequality \(\log x\le x-1\) also gives
\[
L_n\le en\le n^2.
\]
Therefore the target lies in the uniform window
\[
4\le m_{\mathrm{tar},n,\sigma}\le L_n\le n^2.
\]
Choose \(j\ge0\) such that
\[
2^j\le\frac{m_{\mathrm{tar},n,\sigma}}2<2^{j+1},
\qquad
m:=2^j+1.
\]
The lower bound on the target implies \(j\ge1\), so \(m\) is odd. The rounding inequalities give
\[
\frac{m_{\mathrm{tar},n,\sigma}}4\le m\le m_{\mathrm{tar},n,\sigma},
\]
or equivalently,
\[
\frac{L_n}{16384\Lambda_{\mathrm M,n,\sigma}}
\le m
\le\frac{L_n}{4096\Lambda_{\mathrm M,n,\sigma}}.
\]
Since \(m\le n^2\), this polynomial pair belongs to the dictionary by \cref{obj:def:weight-dictionary}.
% lean: CausalSmith.Stat.NoisydoseWeakdesignTransition.polynomial_tuned_dictionary_degree

The same eventual logarithmic-power bound used above places the direct cutoff below the Fourier cutoff. Indeed,
\[
L_n^{d/2}\le L_n^4\le n^{3/8}\le n.
\]
Taking the power \(-1/d<0\) yields
\[
D_n=n^{-1/d}\le L_n^{-1/2}.
\]
Hence, when \(L_n^{-1/2}<\sigma\), the third branch of \cref{obj:def:frontier-rate} applies:
\[
b_{n,\sigma}
=\frac{\Lambda_{\mathrm M,n,\sigma}}{L_n},
\qquad
\rho_{n,\sigma}
=\left(\frac{\Lambda_{\mathrm M,n,\sigma}}{L_n}\right)^\beta.
\]
The lower rounding bound implies
\[
m^{-1}\le16384\,\frac{\Lambda_{\mathrm M,n,\sigma}}{L_n},
\qquad
m^{-\beta}
\le16384^\beta
\left(\frac{\Lambda_{\mathrm M,n,\sigma}}{L_n}\right)^\beta.
\]
Applying the polynomial certificate therefore gives
\[
A_{q_m^{\mathrm M}}^{\mathrm{ref}}
\le C_{\mathrm M}\rho_{n,\sigma},
\qquad
C_{\mathrm M}:=C_{\mathrm{poly}}\,16384^\beta>0.
\]
Taking a common integer threshold \(n_{\mathrm M}\) for the certificate, the target window, and the cutoff comparison proves the polynomial witness uniformly in \(\sigma\).
% lean: CausalSmith.Stat.NoisydoseWeakdesignTransition.inverseheat_dictionary_witness

\item \emph{Selection and the common bound.}
We first convert the reference scores to the public scores of \cref{obj:synth_4}. For every dictionary pair, nonnegativity of the weight and \(|t|^\beta\le1\) on the latent interval give
\[
0\le I_{\kappa+\beta}(q)\le I_\kappa(q)<\infty.
\]
Together with \(I_\kappa(q)>0\), this makes every summand of \(A_q^{\mathrm{ref}}\) nonnegative. The public quantities satisfy
\[
\begin{aligned}
b_q(K)&=p_{\min}c_{\mathrm{lo}}I_\kappa(q),\\
B_q(K)&=\frac{c_{\mathrm{hi}}}{64c_{\mathrm{lo}}}B_q^{\mathrm{ref}},\\
V_q(K)&=\frac{c_{\mathrm{hi}}}{16}V_q^{(16)},\\
\frac{12}{b_q(K)}\sqrt{\frac{V_q(K)}n}
&=\frac{\sqrt{c_{\mathrm{hi}}/16}}{16p_{\min}c_{\mathrm{lo}}}
\frac{12}{b_q^{\mathrm{ref}}}\sqrt{\frac{V_q^{(16)}}n}.
\end{aligned}
\]
The second identity follows from the two bias definitions, and the last follows by factoring the nonnegative constant out of the square root and using \(b_q^{\mathrm{ref}}=I_\kappa(q)/16\). Define
\[
S(K):=\max\left\{1,\frac{c_{\mathrm{hi}}}{64c_{\mathrm{lo}}},
\frac{\sqrt{c_{\mathrm{hi}}/16}}{16p_{\min}c_{\mathrm{lo}}}\right\}\ge1.
\]
Bounding the coefficient of each nonnegative summand by \(S(K)\) yields
\[
\begin{aligned}
A_q(K)
&=\frac{c_{\mathrm{hi}}}{64c_{\mathrm{lo}}}B_q^{\mathrm{ref}}
+\frac{\sqrt{c_{\mathrm{hi}}/16}}{16p_{\min}c_{\mathrm{lo}}}
\frac{12}{b_q^{\mathrm{ref}}}\sqrt{\frac{V_q^{(16)}}n}
+2\sqrt{\frac3n}\\
&\le S(K)\left(B_q^{\mathrm{ref}}
+\frac{12}{b_q^{\mathrm{ref}}}\sqrt{\frac{V_q^{(16)}}n}
+2\sqrt{\frac3n}\right)
=S(K)A_q^{\mathrm{ref}}.
\end{aligned}
\]
% lean: CausalSmith.Stat.NoisydoseWeakdesignTransition.Aq_le_scoreScale_mul_unitAq
Set
\[
C:=S(K)\max\{C_{\mathrm F},C_{\mathrm M}\}>0,
\qquad
n_0:=\max\{3,n_{\mathrm F},n_{\mathrm M}\}.
\]
For \(n\ge n_0\), \(L_n\ge1\). Every branch of the localization scale in \cref{obj:def:frontier-rate} is nonnegative, so
\[
\rho_{n,\sigma}\ge0.
\]
For either family witness \(q\) constructed above, the reference-score bound therefore gives
\[
A_q(K)\le S(K)A_q^{\mathrm{ref}}
\le S(K)\max\{C_{\mathrm F},C_{\mathrm M}\}\rho_{n,\sigma}
=C\rho_{n,\sigma}.
\]
Thus both public witness bounds hold with the same \(C\) and \(n_0\).
% lean: CausalSmith.Stat.NoisydoseWeakdesignTransition.score_le_of_unitScore_le

For completeness, the ordered selection rule gives the required minimizing inequality. Enumerate the dictionary in its fixed order as
\[
N_{\mathrm{dict}}:=|\mathcal D_{n,\sigma}|,
\qquad
q_{(1)},\ldots,q_{(N_{\mathrm{dict}})},
\]
with the constant pair first. Define the successive retained entries by
\[
q_{\mathrm{best},0}:=1,
\qquad
q_{\mathrm{best},i}:=
\begin{cases}
q_{(i)},&A_{q_{(i)}}(K)<A_{q_{\mathrm{best},i-1}}(K),\\
q_{\mathrm{best},i-1},&A_{q_{\mathrm{best},i-1}}(K)\le A_{q_{(i)}}(K),
\end{cases}
\quad 1\le i\le N_{\mathrm{dict}}.
\]
Induction on \(i\) proves
\[
q_{\mathrm{best},i}\in\{1,q_{(1)},\ldots,q_{(i)}\},
\qquad
A_{q_{\mathrm{best},i}}(K)\le A_q(K)
\quad
\text{for every }q\in\{1,q_{(1)},\ldots,q_{(i)}\}.
\]
If the new score is strictly smaller, the new entry satisfies the bound against every preceding entry by transitivity; otherwise the retained entry also bounds the new score. Equal scores retain the earlier entry. At the final index this is the specified first minimizer, and therefore
\[
\widehat q\in\mathcal D_{n,\sigma},
\qquad
A_{\widehat q}(K)\le A_q(K)
\quad(q\in\mathcal D_{n,\sigma}).
\]
% lean: CausalSmith.Stat.NoisydoseWeakdesignTransition.selectedTag_minimizes

Now split at the stated cutoff. If \(\sigma\le L_n^{-1/2}\), the Fourier witness supplies \(k\ge0\) with its pair in the dictionary and
\[
A_{\widehat q}(K)\le A_{q_{2^{-k}}^{\mathrm F}}(K)
\le C\rho_{n,\sigma}.
\]
If \(L_n^{-1/2}<\sigma\), the polynomial witness supplies \(j\ge0\), with \(m=2^j+1\), and
\[
A_{\widehat q}(K)\le A_{q_m^{\mathrm M}}(K)
\le C\rho_{n,\sigma}.
\]
These two cases establish the selected-score bound and both asserted family witnesses with the same constant and threshold.
% lean: CausalSmith.Stat.NoisydoseWeakdesignTransition.dictionary_score_rate
\end{enumerate}
\end{proof}
